\section{Travel of parked mass}\label{sec:travel}

The allocation theorem controls movement between tree vertices.  A second
cost arises when a vertex retains its key but that key acquires a new
anchor in the metric.  We pay this cost using the comparator mass in the
key, normalized by a slowly updated estimate of its size.  The estimate
is essential: keeping its coefficient fixed during filtering allows us
to use conditional expectation without paying the absolute variation of
a posterior.

Throughout this section the tree, label pools, and maps are those of
Section~\ref{sec:partitions}, and the allocations are supplied by
Theorem~\ref{thm:allocation}.  A state at time $t$ means the state after
all updates for that request.  Write $\mu_t$ for the true-metric posterior
counting measure at this state.  All multiplicative constants below are
absolute.  Additive constants may depend on the fixed finite metric,
$k$, the scales, and the fixed label pools, but not on the input law or
its horizon.

\subsection{Parking and valid representatives}

For every tree vertex, including the root and leaves, define its parked
mass by
\begin{equation}\label{eq:park-definition}
 p_{t,v}^{\circ}=q_{t,v}-\sum_{i\text{ child of }v}q_{t,i}.
\end{equation}
The empty sum at a leaf is zero.  Allocation feasibility makes these
numbers nonnegative.  Telescoping within a subtree gives
\begin{equation}\label{eq:park-subtree-total}
 \sum_{w\text{ descendant of }v}p_{t,w}^{\circ}=q_{t,v},
 \qquad \sum_v p_{t,v}^{\circ}=k,
\end{equation}
where a vertex counts as its own descendant.

\begin{lemma}[Parked variation and witnesses]\label{lem:park-variation}
At every step,
\begin{equation}\label{eq:park-variation}
 \sum_v r_{|v|}\lvert\Delta p_v^{\circ}\rvert
 \le \left(1+\frac1\tau\right)
       \sum_{v\ne\mathrm{root}}r_{|v|-1}\lvert\Delta q_v\rvert.
\end{equation}
If $q_{t,v}>0$, there is a point of $Y$ whose current key string has
prefix $v$.  Such a point is called a validity witness for $v$.  If
$|v|=j\ge1$, every validity witness is within $20r_j$ of the anchor of
the last key of $v$.
\end{lemma}

\begin{proof}
Taking differences in \eqref{eq:park-definition} and using the triangle
inequality gives
\begin{align*}
 \sum_v r_{|v|}\lvert\Delta p_v^{\circ}\rvert
 &\le \sum_v r_{|v|}\lvert\Delta q_v\rvert
     +\sum_v r_{|v|}\sum_{i\text{ child of }v}
                                  \lvert\Delta q_i\rvert\\
 &=\frac1\tau\sum_{v\ne\mathrm{root}}
                 r_{|v|-1}\lvert\Delta q_v\rvert
   +\sum_{i\ne\mathrm{root}}r_{|i|-1}\lvert\Delta q_i\rvert.
\end{align*}
The root term in the first sum vanishes because its allocation is always
$k$.  This proves \eqref{eq:park-variation}, including the root's parked
mass.  The bound $q_{t,v}\le C_0m_{t,v}$ implies $m_{t,v}>0$ whenever
$q_{t,v}>0$.  Since $Y$ is finite, some point mapped below $v$ then has
positive posterior counting mass.  It is a witness.  The anchor-distance
assertion is part of Lemma~\ref{lem:fixed-alphabet}: a level key is a
heavy ball of radius at most $20r_j$, a smaller roster ball, or a
singleton.
\end{proof}

Represent a park at a nonroot vertex by its last key's anchor.  Represent
the root by one fixed point of $Y$.  A balanced integer rounding of
$p_{t,v}^{\circ}$ can be positive only if $p_{t,v}^{\circ}>0$, since the
floor and ceiling of zero are both zero.  Thus the same witness statement
will apply to occupied integer parks in Section~\ref{sec:rounding}.

Let $\Lambda_j^H$ be the fixed pool of heavy-label slots at level $j$.
For $\lambda\in\Lambda_j^H$, let $C_{t,j,\lambda}\subseteq Y$ be the
points currently assigned that level key, and let $a_{t,j,\lambda}$ be
its anchor.  An absent label has empty region and an arbitrary fixed
placeholder anchor.  Heavy precedence and disjointness mean that, when
present, $C_{t,j,\lambda}$ is the entire heavy ball carrying the label.
Set
\begin{equation}\label{eq:label-park-size}
 \begin{aligned}
 b_{t,j,\lambda}
   &=\sum_{\substack{|v|=j\\\text{last key of }v=\lambda}}
                                              p_{t,v}^{\circ},\\
 M_{t,j,\lambda}
   &=\sum_{\substack{|v|=j\\\text{last key of }v=\lambda}}m_{t,v}
     =\mu_t(C_{t,j,\lambda}).
 \end{aligned}
\end{equation}
These definitions use the fixed slots even when they are absent, in
which case both quantities are zero.  In particular,
\begin{equation}\label{eq:label-park-domination}
 0\le b_{t,j,\lambda}\le C_0 M_{t,j,\lambda},\qquad
 \sum_{j,\lambda\in\Lambda_j^H}r_j\lvert\Delta b_{j,\lambda}\rvert
 \le\sum_v r_{|v|}\lvert\Delta p_v^{\circ}\rvert.
\end{equation}
The second inequality follows by grouping fixed vertex coordinates,
then using the triangle inequality.  It also holds across label
disappearance and reuse.

Let $E_{t,j,\lambda}$ be the event that the label is present at both
successive states and its anchor changes.  By
Lemma~\ref{lem:fixed-alphabet}, this is an explicit reuse of a heavy
label at an intersecting new heavy ball.  In particular,
\begin{equation}\label{eq:retained-anchor-distance}
 \one_{E_{t,j,\lambda}}
 d(a_{t-1,j,\lambda},a_{t,j,\lambda})
 \le 40r_j\one_{E_{t,j,\lambda}}.
\end{equation}
The quantity to be bounded is
\begin{equation}\label{eq:anchor-travel-definition}
 T_{\rm anchor}
   =\E\sum_{t,j,\lambda\in\Lambda_j^H}
                  r_j b_{t,j,\lambda}\one_{E_{t,j,\lambda}}.
\end{equation}

\begin{proposition}[Retained-anchor travel]\label{prop:anchor-travel}
There is an absolute constant $C$ and a finite constant $D_{\rm anchor}$,
independent of the input law and horizon, such that
\begin{equation}\label{eq:anchor-travel-bound}
 T_{\rm anchor}\le C\bigl(V+\ell P+\ell Q\bigr)+D_{\rm anchor}.
\end{equation}
Consequently the expected cost of moving the new parked masses with
their retained anchors is at most $40T_{\rm anchor}$.  No multiplicative
constant in this statement depends on the number of levels or on $k$.
\end{proposition}

The proof uses a size tracker and a pointwise estimate for its old
coefficients.  We establish both before analyzing the travel potential.

\subsection{An affine size tracker}

Put $\rho=6/5$.  For each fixed heavy-label slot maintain a reference
$\widehat m_{t,j,\lambda}\in[1,k+1]$ for
$s_{t,j,\lambda}=1+M_{t,j,\lambda}$.  Initially set the reference equal
to $s_{0,j,\lambda}$.  At each subsequent step, after computing the new
membership and posterior size, keep the old reference if
\begin{equation}\label{eq:key-size-comparison}
 \rho^{-1}\widehat m_{t-1,j,\lambda}
       \le s_{t,j,\lambda}
       \le\rho\widehat m_{t-1,j,\lambda};
\end{equation}
otherwise refresh it to $s_{t,j,\lambda}$.  Thus every completed state
satisfies the same comparison with its current reference.  The tracker
continues to run while a slot is absent; recycling a slot does not
restart it.

\begin{lemma}[Size-refresh budget]\label{lem:key-size-refresh}
If $R_{t,j,\lambda}$ denotes a refresh, then
\begin{equation}\label{eq:key-size-refresh-budget}
 \E\sum_{t,j,\lambda\in\Lambda_j^H}
 r_j\bigl(\widehat m_{t-1,j,\lambda}
            +\widehat m_{t,j,\lambda}\bigr)\one_{R_{t,j,\lambda}}
 \le C\ell P+D_{\rm size},
\end{equation}
where $C$ is absolute and $D_{\rm size}$ depends only on the fixed
metric, scales, $k$, and label pools.
\end{lemma}

\begin{proof}
The payment is a tangent gap, equivalently the Bregman divergence
for the convex function $s\log s$ \cite{Bregman67}.
For $u,s\in[1,k+1]$, the tangent at $u$ to $f(s)=-s\log s$ is
\begin{equation}\label{eq:entropy-tangent}
 T_u(s)=f(u)+f'(u)(s-u)=u-(1+\log u)s.
\end{equation}
It is affine in $s$, with slope of absolute value at most $\ell$.  A
refresh $u\to s$ releases
\begin{equation}\label{eq:entropy-refresh-credit}
 T_u(s)-T_s(s)=s\log(s/u)-s+u
       \ge\frac{(s-u)^2}{2\max(s,u)}.
\end{equation}
For completeness, if $s\ge u$ the left side equals
$\int_u^s(s-z)z^{-1}\,dz$, and if $s\le u$ it equals
$\int_s^u(z-s)z^{-1}\,dz$.  Bounding the denominator by
$\max(s,u)$ gives the inequality.  At failure of
\eqref{eq:key-size-comparison},
$|s-u|\ge(1-\rho^{-1})\max(s,u)$, so
\begin{equation}\label{eq:entropy-relative-credit}
 T_u(s)-T_s(s)\ge\frac{(1-\rho^{-1})^2}{4}(s+u)
                         =\frac{s+u}{144}.
\end{equation}

We check the input budget for the weighted sum of these tangents.  Let
$N_{t,j,\lambda}^{\rm key}$ be the actual number of hidden members whose
level-$j$ key at time $t$ is $\lambda$.  Filtering replaces
$\E[N_{t-1,j,\lambda}^{\rm key}\mid\mathcal F_{t-1}]$ by
$\E[N_{t-1,j,\lambda}^{\rm key}\mid\mathcal F_t]$, using the old map
and old reference.  Its signed contribution has conditional expectation
zero.  Drift then replaces this old hidden count by the new count,
with both conditioned on $\mathcal F_t$.

If a member's two mapped leaves first differ at level $j_0$, the
radius sum of its possible level-key disagreements is at most
\begin{equation}\label{eq:key-disagreement-sum}
 \sum_{j\ge j_0}r_j
       \le\frac{r_{j_0-1}}{\tau-1}
       =\frac{d_T(\text{old leaf},\text{new leaf})}{\tau-1}.
\end{equation}
At one level a member changes at most two heavy-slot indicators, one
for the departing key and one for the arriving key.  Applying this
pathwise to the coupled hidden trajectory and then taking conditional
expectations gives
\begin{equation}\label{eq:key-count-drift}
 \E\sum_{t,j,\lambda\in\Lambda_j^H}r_j
 \left|M_{t,j,\lambda}
       -\E[N_{t-1,j,\lambda}^{\rm key}\mid\mathcal F_t]\right|
 \le\frac{2}{\tau-1}P\le2P.
\end{equation}
Thus the tangents' expected total signed input increase is at most
$2\ell P$.

The tangents need not be nonnegative.  Their endpoint values nevertheless
have a uniform bound:
\[
 |T_u(s)|\le(k+1)(1+\ell)
       \qquad(1\le u,s\le k+1).
\]
Write $L_H=\sum_j r_j|\Lambda_j^H|$, a finite quantity fixed before
the input law or horizon is chosen.  The absolute value of the summed
tangent potential is at most $(k+1)(1+\ell)L_H$ at both endpoints.
Telescoping its signed changes over every step, and using
\eqref{eq:entropy-relative-credit}, proves
\eqref{eq:key-size-refresh-budget}, for example with
$C=288$ and
$D_{\rm size}=288(k+1)(1+\ell)L_H$.  Filtering is included on all steps
before the decision to refresh; it is not restricted to the adaptively
chosen refresh events.
\end{proof}

\subsection{A pointwise bound on old parked coefficients}

At a completed state, \eqref{eq:label-park-domination} and the size
comparison give the individual bound
\begin{equation}\label{eq:individual-park-coefficient}
 \frac{b_{j,\lambda}}{\widehat m_{j,\lambda}}
 \le\rho C_0\frac{M_{j,\lambda}}{1+M_{j,\lambda}}
 \le\rho C_0.
\end{equation}
Summing this bound over all levels would introduce the height of the
tree.  The useful additional constraint is that parks at different depths
draw on the same subtree allocations.  All finer parks near a point can
therefore be charged to comparator mass in a slightly larger ball at a
coarser scale.  Meanwhile a heavy ball whose anchor is near that point
contains a smaller ball and bounds the normalization from below.
Comparing these two ball masses will give a logarithmic coefficient sum.

\begin{lemma}[Pointwise parked coefficients]\label{lem:park-coefficients}
Fix any completed fractional state and any $y\in Y$.  At each level
whose heavy-center set has a center within $7r_j$ of $y$, let
$\lambda_j$ be its label.  Such a center is unique.  Then
\begin{equation}\label{eq:pointwise-park-coefficients}
 \sum_{\substack{j:\ d(y,a_{j,\lambda_j})\le7r_j}}
                \frac{b_{j,\lambda_j}}{\widehat m_{j,\lambda_j}}
 \le C\ell,
\end{equation}
with an absolute $C$, independent of the state and the number of levels.
\end{lemma}

\begin{proof}
Call the levels in the sum selected, and abbreviate their numerator to
$b_j$.  Distinct heavy centers at one level are separated by more than
$100r_j$, proving uniqueness.  Their radii are at least $16r_j$;
therefore the selected heavy ball contains $B(y,r_j)$.  Heavy precedence
and the size comparison imply
\begin{equation}\label{eq:park-denominator-lower}
 \widehat m_{j,\lambda_j}
       \ge\rho^{-1}\bigl(1+\mu(B(y,r_j))\bigr).
\end{equation}
Here $\mu$ is the posterior counting measure of the fixed state.

We also have, for every level $j$, the cumulative capacity estimate
\begin{equation}\label{eq:cumulative-park-mass}
 \sum_{\substack{i\ge j\\i\text{ selected}}}b_i
                      \le C_0\mu(B(y,67r_j)).
\end{equation}
To prove it, consider each positive park counted on the left.  At its
depth $i$, a validity witness lies within $20r_i$ of its anchor, hence
within $27r_i$ of $y$.  Group all these parks by their depth-$j$
ancestor $v$.  Their sum in each group is at most $q_v$ by
\eqref{eq:park-subtree-total}.  The group has at least one witness $z$
within $27r_i\le27r_j$ of $y$.  Every point with prefix $v$ shares its
level-$j$ cell with $z$, and that cell has diameter at most $40r_j$.
All such points therefore belong to $B(y,67r_j)$.  Regions of distinct
depth-$j$ prefixes are disjoint, so their posterior counts sum to at
most $\mu(B(y,67r_j))$.  Summing $q_v\le C_0m_v$ proves
\eqref{eq:cumulative-park-mass}.

For an explicit scale summation put, at every level,
\[
 A_j=1+\mu(B(y,67r_j)),\qquad B_j=1+\mu(B(y,r_j)).
\]
Both are between $1$ and $k+1$.  Extend $B_j=k+1$ for $j\le0$.
Since $\tau\ge2$, the fixed offset $p=7$ satisfies $\tau^p\ge67$,
and hence $A_j\le B_{j-p}$.  It follows that
\begin{equation}\label{eq:park-log-telescope}
 \sum_{j=1}^J\log(A_j/B_j)
 \le\sum_{j=1}^J\bigl(\log B_{j-p}-\log B_j\bigr)
 \le p\log(k+1).
\end{equation}
The final inequality cancels the common middle terms; at most $p$
uncancelled positive terms remain, each at most $\log(k+1)$, and the
remaining terms are nonpositive.  This includes $J<p$.

There are at most $p\log(k+1)/\log2$ levels with $A_j>2B_j$.
Their contribution is $O(\ell)$ by
\eqref{eq:individual-park-coefficient}.  Partition the remaining
selected levels into groups
\[
 G_a=\{j:\ j\text{ selected},\ A_j\le2B_j,
                         \ 2^a\le A_j<2^{a+1}\},
 \qquad 0\le a\le\lfloor\log_2(k+1)\rfloor.
\]
For a nonempty group let $j_a$ be its smallest, hence coarsest, index.
By \eqref{eq:park-denominator-lower}, every denominator in this group
is at least $2^{a-1}/\rho$.  By
\eqref{eq:cumulative-park-mass},
\[
 \sum_{j\in G_a}b_j
  \le C_0\mu(B(y,67r_{j_a}))
  =C_0(A_{j_a}-1)\le C_0 2^{a+1}.
\]
Its total coefficient is consequently at most $4\rho C_0$.  There are
at most $1+\log_2(k+1)$ groups, proving
\eqref{eq:pointwise-park-coefficients}.  All these statements hold for
each individual state and point, with no averaging and no conditioning
on a future request.
\end{proof}

\subsection{Paying retained-anchor motion}

\begin{proof}[Proof of Proposition~\ref{prop:anchor-travel}]
Use the clipped ramp
\[
 \kappa(u)=
 \begin{cases}
  0,&u\le4,\\
  (u-4)/3,&4<u<7,\\
  1,&u\ge7.
 \end{cases}
\]
For a completed state define
\begin{equation}\label{eq:anchor-potential}
 \Phi_t=\sum_{j,\lambda\in\Lambda_j^H}
    r_j\frac{b_{t,j,\lambda}}{\widehat m_{t,j,\lambda}}
    \int_{C_{t,j,\lambda}}
           \kappa\!\left(\frac{d(a_{t,j,\lambda},z)}{r_j}\right)
                                                   \,d\mu_t(z).
\end{equation}
This is the posterior expectation of the sum of the same terms over
the hidden members assigned to each key.  The ramp measures mass away
from the anchor.  When a heavy label is reused at a new request, the
concentrated mass is far from its old anchor and close to its new one,
so the anchor change decreases the potential.  Dividing by the held
key size makes this decrease comparable to $r_jb_{t,j,\lambda}$,
the travel charge, rather than to its product with comparator mass.

The potential is nonnegative.  Moreover,
at a completed state the normalized integral is at most $\rho$, so
\begin{equation}\label{eq:anchor-potential-endpoints}
 0\le\Phi_t\le\rho\sum_{j,\lambda\in\Lambda_j^H}r_jb_{t,j,\lambda}
                         \le\rho k r_0.
\end{equation}
The last inequality follows from \eqref{eq:park-subtree-total}, since
each park occurs at just one depth and $r_j\le r_0$.

We now decompose the change of \eqref{eq:anchor-potential} in a specified
order.  This is an accounting order: the current partition and allocation
are already available from the fractional history.  Old parameters are
held until their indicated substep.

\paragraph{Filtering with old membership and all old parameters.}
In the hidden-member expression for $\Phi_{t-1}$ retain the old locations,
old key map, old anchors, old parks, and old references.  Replace
conditioning on $\mathcal F_{t-1}$ by conditioning on $\mathcal F_t$.
The coefficients and the old map are $\mathcal F_{t-1}$-measurable.
Thus the tower property makes this signed potential change have
conditional expectation zero.  We perform this comparison on every
step, before any refresh or anchor event is selected.

\paragraph{Drift to new membership and new locations.}
Use the coupled hidden trajectory, now conditioned on $\mathcal F_t$,
but continue to hold all anchors, parks, and size references at their
old values.  For a member whose level key changes, its total possible
positive change among heavy slots is at most $C r_j$, by
\eqref{eq:individual-park-coefficient} and $0\le\kappa\le1$.
Summing \eqref{eq:key-disagreement-sum} over members and steps pays
these terms by $CP$ in expectation.  This includes changes to heavy
membership caused by partition edits, regardless of physical movement.

For unchanged membership in a heavy key, the same old anchor appears
in both ramp terms.  If the old location is $y$ and the new one is $z$,
then
\begin{equation}\label{eq:anchor-unchanged-key-drift}
 r_j\frac{b^{\rm old}_{j,\lambda}}{\widehat m^{\rm old}_{j,\lambda}}
 \left[\kappa\!\left(\frac{d(a^{\rm old}_{j,\lambda},z)}{r_j}\right)
       -\kappa\!\left(\frac{d(a^{\rm old}_{j,\lambda},y)}{r_j}\right)
 \right]_+
 \le\frac{d(y,z)}3
       \frac{b^{\rm old}_{j,\lambda}}{\widehat m^{\rm old}_{j,\lambda}}
       \one_{\{d(a^{\rm old}_{j,\lambda},y)<7r_j\}}.
\end{equation}
Indeed $\kappa$ is $1/3$-Lipschitz and cannot increase from a point where
it already equals one.  Lemma~\ref{lem:park-coefficients} bounds the
sum of the coefficients on the right by $C\ell$, for every old point
$y$ and the old fractional state.  This pointwise bound remains valid
when that $y$ is sampled in the coupling conditional on the new history;
we require no independence between the old mover and the new request.
Only the designated comparator mover changes physical position on this
step.  Summing \eqref{eq:anchor-unchanged-key-drift} therefore costs at
most $C\ell e_x$ conditionally, and at most $C\ell Q$ in total
expectation.  The entire drift budget is $C(P+\ell Q)$.

\paragraph{Refreshing the size references.}
At this point each integral uses the new key membership and new
posterior locations, but still the old anchor.  For one label write
this integral as $I$, its new size as $M$, and $s=1+M$.  Always
$0\le I\le M$.  At a refresh from $u=\widehat m^{\rm old}$ to $s$ the
change is
\[
 r_j b^{\rm old}I\left(\frac1s-\frac1u\right).
\]
If $s\ge u$ this is nonpositive.  If $s<u$, its positive part is at
most $r_jb^{\rm old}I/s\le r_jb^{\rm old}$.  The old completed state
gives $b^{\rm old}\le\rho C_0 u$ by
\eqref{eq:individual-park-coefficient}.  Hence all positive refresh
jumps are at most a constant times the charge in
\eqref{eq:key-size-refresh-budget}.  This argument does not assume that
the new integral is controlled by the old denominator.

\paragraph{Updating the parked amounts.}
The new size references now satisfy
$\widehat m^{\rm new}\ge(1+M)/\rho$, whether refreshed or retained.
Consequently $I/\widehat m^{\rm new}\le\rho$.  Replacing $b^{\rm old}$
by $b^{\rm new}$ changes the potential by at most
$\rho r_j|\Delta b|$ in absolute value.  Equations
\eqref{eq:label-park-domination} and \eqref{eq:park-variation} give
the expected total budget
\begin{equation}\label{eq:anchor-park-update-budget}
 \rho\E\sum_{t,j,\lambda\in\Lambda_j^H}r_j|\Delta b_{j,\lambda}|
                      \le\rho(1+\tau^{-1})V.
\end{equation}

\paragraph{Updating anchors and collecting the travel payment.}
An absent final key has zero integral and zero new parked mass, so
choosing its placeholder anchor costs nothing.  A newly present key
had $b^{\rm old}=0$: the fresh-label rule ensures that it was absent
from the preceding map.  Its anchor initialization can increase the
potential by at most
\[
 \rho r_j b^{\rm new}=\rho r_j|\Delta b|,
\]
which is another copy of the variation budget
\eqref{eq:anchor-park-update-budget}.  A retained key with unchanged
anchor contributes zero.

It remains to consider a retained moving heavy label at an update for
the request $x$, at radius $r=r_j$.  Its new anchor is $x$, its old
anchor $a$ satisfies $d(a,x)>8r$, and its new key region is the entire
new heavy ball $C$.  Put
\[
 H=\mu_t(B(x,\gamma_H r)).
\]
The served comparator contains $x$, so $H\ge1$.  The new heavy ball
contains this inner ball and is contained in $B(x,D_Hr)$.  The heavy
condition \eqref{eq:heavy-condition} therefore gives
\begin{equation}\label{eq:anchor-heavy-mass-comparison}
 H\le M^{\rm new}_{j,\lambda}=\mu_t(C)\le(1+\delta_H)H,
 \qquad
 \widehat m^{\rm new}_{j,\lambda}
       \le\rho(1+M^{\rm new}_{j,\lambda})
       \le\rho(2+\delta_H)H.
\end{equation}
The last inequality is exactly where the inner mass $H\ge1$ is needed;
it accounts for the added one in the normalization.

On $B(x,\gamma_Hr)$ the old ramp is one because
$d(a,z)>(8-\gamma_H)r>7r$, and the new ramp is zero because
$d(x,z)\le\gamma_Hr<4r$.  Outside this inner ball but inside $C$, the
ramp can increase by at most one, and that region has mass at most
$\delta_H H$.  Hence changing $a$ to $x$ changes the unnormalized
integral by at most $-(1-\delta_H)H$.  With the new park and new
reference already in place, \eqref{eq:anchor-heavy-mass-comparison}
shows that the potential decreases by at least
\begin{equation}\label{eq:retained-anchor-drop}
 \frac{1-\delta_H}{\rho(2+\delta_H)}\,r b^{\rm new}_{j,\lambda}
                         \ge\frac14 r b^{\rm new}_{j,\lambda}.
\end{equation}
Here we used the absolute choices $\gamma_H<1/100$ and
$\delta_H<1/100$.  Thus this payment is uniform even when the retained
park remains at exactly the same tree vertex.

Combining the five substeps, summing signed changes over all times, and
taking expectations yields
\[
 \frac14 T_{\rm anchor}
 \le \E\Phi_0-\E\Phi_T
       +C(P+\ell Q)+C\ell P+C D_{\rm size}+CV.
\]
The filtering contributions have vanished by conditional expectation;
the size-refresh estimate pays their later adaptively selected
consequences.  Use \eqref{eq:anchor-potential-endpoints} and $\ell\ge1$
to obtain \eqref{eq:anchor-travel-bound}, with an additive constant
bounded by $C(D_{\rm size}+kr_0)$.  Finally,
\eqref{eq:retained-anchor-distance} proves the asserted bound for actual
anchor displacement.

The only scale summations used were the geometric sum
\eqref{eq:key-disagreement-sum} and the fixed-offset and dyadic sums in
Lemma~\ref{lem:park-coefficients}.  Their constants are independent of
the height.  The finite label pools enter only the endpoint constant
$D_{\rm size}$, and their persistent tracker states ensure that no
initialization allowance is charged again when a label is recycled.
\end{proof}
